\begin{figure}[h]
	\centering
	\tikzstyle{block} = [rectangle, draw, minimum width = 2cm, minimum height = 2cm]  % bloc
	%
	%
	\tikzsetnextfilename{fig_interaction_port}
	\begin{tikzpicture} [node distance = 2cm]
		\node (O) at (0,0) {};
		\node[block, left of = O, align=center] (robot) {Système 1};
		\node[block, right of = O, align=center] (env) {Système 2};
		%
		\draw[>=Latex, <->] ($(robot.east) + (0,0.5)$) -- ($(env.west) + (0,0.5)$)
			node[midway, above] {$\vect{f}$};
		\draw[>=Latex, <->] ($(robot.east) - (0,0.5)$) -- ($(env.west) - (0,0.5)$)
		node[midway, above] {$\vect{v}$};
		%
	\end{tikzpicture}
	\caption{Exemple de ports d'interaction entre deux systèmes physiques}
	\label{fig_physical_interaction}
\end{figure}